\infographicchapter{Terahertz Imaging}{thz.png}
\label{ch:infographic-terahertz-imaging}

\section{Overview}

Terahertz (THz) imaging---sometimes informally called ``T-rays''---is a research-stage, radiation-free medical scan that uses electromagnetic waves sitting right between microwaves and infrared light (typically 0.1 to 3.0~terahertz within the broader 0.1--10~THz band). Because terahertz waves carry extremely tiny amounts of energy, they cannot ionize molecules or harm DNA, making them naturally safe for repeated exposure. Their biggest strength is how strongly they react to water and gentle molecular motions in tissue. Because cancerous spots, swollen tissue, and burns often hold different amounts of water than healthy tissue around them, terahertz waves can highlight those differences without any injected dyes.

Scientists test both steady-beam and ultra-short pulsed terahertz scanners (called terahertz pulsed imaging). By measuring how much of the wave bounces back and how slightly it slows down, computers can map surface water content and tissue structure. Historically, building compact equipment for this ``terahertz gap'' has been challenging because ordinary electronics are too slow for these frequencies while standard optical lasers carry too much energy; however, newer lab prototypes are adding compact solid-state chips, quantum cascade lasers, smart sensors, and machine learning to make scans faster and clearer.

Today, terahertz imaging is strictly in the laboratory and early prototype research stage---no terahertz scanner is approved for routine hospital diagnosis. Researchers are focusing mainly on surface checks, such as mapping the edges of skin cancers (like basal cell carcinoma), measuring eye-surface (corneal) swelling, checking burn depth, and scanning freshly removed breast lumpectomy tissue in the lab to see if cancer reaches the edge. The main hurdles are that body water blocks terahertz waves within just a millimeter or two, harmless dense tissue can look similar to cancer, and the hardware is still bulky and costly.

\section{Underlying Physical Principles}
\begin{itemize}
  \item \textbf{Ultra-Short Pulses or Steady Terahertz Beams:} Pulsed terahertz scanners use ultrafast lasers hitting special tiny antennas or crystals to send out sub-picosecond bursts of terahertz waves, while newer solid-state chips create steady continuous beams.
  \item \textbf{Extreme Sensitivity to Water:} Low-energy terahertz waves are completely non-ionizing and interact strongly with free and bound water molecules as well as gentle vibrations between large biological molecules---creating natural contrast for tumor swelling, burn depth, and corneal hydration.
  \item \textbf{Surface-Only Reflection Scanning:} While terahertz waves can pass fairly easily through dry, non-metallic materials (such as dry dressings or outer tooth enamel) that block visible light, liquid water absorbs them very strongly. Because of that absorption, the waves travel only about 0.1 to 2~millimeters into moist tissue, so medical prototypes work in reflection mode, bouncing waves off the skin, eye surface, or freshly removed surgical tissue.
  \item \textbf{Precise Timing and Strength Detection:} Sensitive detectors synchronized with the laser pulse measure both the strength and exact arrival timing of the returning terahertz wave while ignoring everyday room heat.
  \item \textbf{Mapping Water and Tissue Properties:} Comparing the bounced wave from tissue against a reference mirror lets computers calculate how much the tissue slowed and absorbed the wave, feeding hydration maps and research classifiers.
\end{itemize}

\section{Imaging Mechanism}
\begin{enumerate}
  \item \textbf{Non-Contact Alignment and Calibration:} For research skin, burn, or eye scans, the area is held steadily in front of a non-contact optical head or a clean quartz window (freshly removed surgical specimens rest flat on a lab stage). Staff calibrate against a mirror and sometimes flush the beam path with dry nitrogen gas so room humidity does not weaken the signal.
  \item \textbf{Scanning the Terahertz Beam:} An ultrafast pump laser triggers a tiny emitter to send gentle terahertz pulses onto the surface while mirrors or motorized stages scan the beam across the target area.
  \item \textbf{Capturing the Reflected Wave:} Returning terahertz waves focus onto a synchronized detector that records the exact wave strength and timing as digital data.
  \item \textbf{Building Hydration and Tissue Maps:} Software compares the tissue signals against the mirror baseline, extracts surface water and refractive properties, and runs research algorithms to highlight unusual zones.
  \item \textbf{Displaying Research Maps:} The research computer shows 2D color maps of surface reflection, hydration, or experimental tumor-edge scores so scientists can compare them later with official lab pathology or eye OCT scans.
\end{enumerate}

\section{Key Technical Parameters}
\noindent
\begin{tabularx}{\textwidth}{@{}l X@{}}
\toprule
\textbf{Parameter} & \textbf{Typical Value / Description} \\
\midrule
Physical Principle & Bouncing non-ionizing terahertz waves off tissue to measure surface water content and gentle molecular vibrations \\
Energy / Signal Source & Ultrafast laser-driven photoconductive antennas (pulsed THz) or solid-state continuous-wave THz emitters \\
Typical Frequency / Energy & 0.1--3.0~THz (ultra-low photon energy, far below anything that could ionize DNA) \\
Spatial Resolution & About 0.1--1.0~mm across the surface; depth strictly limited to 0.1--2.0~mm by water absorption \\
Typical Scan Duration & Seconds to a few minutes per surface or lab specimen sweep (15--30~minutes total research visit) \\
Contrast Agent & None routinely needed (relies on natural water differences; clearing agents studied only in labs) \\
Ionizing Radiation & None (uses safe, non-ionizing terahertz waves) \\
\bottomrule
\end{tabularx}

\section{Clinical Usages}

\subsection{Primary Usages}
\begin{itemize}
  \item \textbf{Checking Fresh Breast Lumpectomy Edges (Lab Research):} After breast-conserving surgery, standard pathology can take around 10~days and roughly 15--20\% of patients need a second operation for unclear margins. Researchers scan freshly removed breast cancer specimens with terahertz systems---including compact handheld lab prototypes reaching 87--88\% accuracy in early excised-tissue tests---to see if higher-water cancer tissue stands out rapidly from surrounding fat.
  \item \textbf{Mapping Skin Cancer Borders (Early Research):} Because terahertz waves see just the top millimeters of skin, reflection prototypes are studied to outline basal cell carcinoma and melanoma borders before skin surgery.
  \item \textbf{Measuring Eye-Surface Swelling (Translational Research):} Non-contact terahertz sensors are being tested to measure water content in the clear front window of the eye (cornea) for conditions like Fuchs dystrophy and corneal swelling.
  \item \textbf{Assessing Burn Depth (Preclinical / Pilot Research):} Early studies test whether terahertz sensitivity to fluid buildup can help tell shallow burns from deep burns that need skin grafts.
\end{itemize}

\subsection{Secondary Usages}
\begin{itemize}
  \item \textbf{Early Tooth-Decay Research (Lab Studies):} Scientists explore terahertz scans on extracted teeth to spot early enamel thinning and cavities without dental X-rays.
  \item \textbf{Digestive, Brain, and Scar-Healing Studies (Lab Research):} Proof-of-concept studies test freshly removed colon, stomach, and brain tumor samples (seeking dye-free tumor borders without injected PET tracers or fluorescent agents), while early skin studies track wound healing and distinguish raised hypertrophic scars from normal scars over months of recovery.
  \item \textbf{Biosensing and Pill Quality Checks (Lab / Industry):} Terahertz sensors are studied to detect tiny biological molecules and inspect pill coatings in pharmaceutical factories.
\end{itemize}

\section{Advantages}
\begin{itemize}
  \item \textbf{Completely Non-Ionizing and Safe:} Terahertz waves carry far too little energy to harm DNA, making them safe for repeated surface measurements.
  \item \textbf{Highly Sensitive to Water Changes:} Strong interaction with water provides natural, dye-free contrast for tumor swelling, corneal hydration, and burn depth.
  \item \textbf{No Touching, Squeezing, or Dyes Needed:} Reflection systems can check skin, eyes, or fresh surgical tissue without contrast injections, harsh stains, or painful pressure.
  \item \textbf{Rich Quantitative Data:} Pulsed terahertz systems measure both wave strength and timing across many frequencies at once, giving detailed surface hydration metrics.
\end{itemize}

\section{Disadvantages}
\begin{itemize}
  \item \textbf{Stopped Quickly by Body Water:} Because water absorbs terahertz waves intensely, they only reach a fraction of a millimeter to about 2~millimeters deep---never deep internal organs.
  \item \textbf{Less Sharp Than Optical Microscopes:} Sub-millimeter terahertz waves cannot see individual cells the way OCT, confocal microscopy, or standard lab pathology can.
  \item \textbf{False Alarms from Normal Moist Tissue:} Dense harmless breast tissue, scars, and everyday inflammation also hold lots of water, so they can look very similar to cancer.
  \item \textbf{Bulky, Slow, and Expensive Gear:} Pulsed terahertz lasers and optical delay lines are still large, costly, and too slow for busy everyday operating rooms.
  \item \textbf{Strictly Investigational:} No terahertz device is approved for routine clinical diagnosis; everyday use awaits smaller chips, standard rules, and large clinical trials.
\end{itemize}

\section{Contraindications}

\subsection{Absolute}
\begin{itemize}
  \item \textbf{No Standalone Medical Use and Pump-Laser Safety:} Because no terahertz scanner is clinically approved, it must never be used alone to make medical or surgical decisions outside IRB-approved research. Exposure to unenclosed ultrafast pump lasers without eye protection is strictly forbidden.
  \item \textbf{Surface Heating Limits:} Power levels above international skin or eye thermal limits (ICNIRP / IEEE~C95.1) are strictly contraindicated, especially on the cornea or fragile burn tissue.
\end{itemize}

\subsection{Relative}
\begin{itemize}
  \item \textbf{Wet Bandages, Creams, or Deep Targets:} Ointments, damp dressings, active bleeding, or lumps deeper than 1 to 2~millimeters block terahertz waves completely and prevent a useful reading.
\end{itemize}

\section{Patient Safety Guidelines}
\begin{itemize}
  \item \textbf{Protecting Skin and Eyes from Warmth:} Because surface water absorbs terahertz energy in the top millimeters, prototypes keep power and scan time strictly within ICNIRP and IEEE~C95.1 limits so skin and eyes never warm up.
  \item \textbf{Pump-Laser Eye Safety:} Pulsed terahertz scanners use hidden near-infrared lasers inside to create the terahertz beam; strict enclosures, safety switches, and protective eyewear (ANSI~Z136 / IEC~60825) keep eyes completely safe.
  \item \textbf{Ethics Oversight and Standard Care First:} All human or fresh-tissue studies run under approved IRB protocols with written consent, and terahertz results never replace standard pathology or proven medical imaging.
  \item \textbf{Electrical and Surface Hygiene:} Prototypes pass medical electrical safety checks, any quartz contact window is sanitized between participants, and fresh surgical specimens follow hospital biohazard rules.
\end{itemize}

\section{Observing for Abnormal Patterns}
\begin{itemize}
  \item \textbf{Higher Water Signals in Cancer Edges:} On removed breast and skin cancer research samples, tumors typically bounce back more terahertz energy and absorb more signal than surrounding fat because cancer cells hold more water and pack tightly.
  \item \textbf{Burn Depth and Swelling Shifts:} In wound research, deeper burns show different terahertz reflection patterns than intact skin due to fluid buildup and heat-damaged blood vessels.
  \item \textbf{Swollen Corneas:} In eye research, higher terahertz reflection tracks extra water buildup in conditions like Fuchs dystrophy.
  \item \textbf{Overlap with Harmless Dense Tissue:} Normal dense breast glands, scars, and simple inflammation also hold extra water, often looking very similar to cancer without advanced computer classifiers.
  \item \textbf{Surface Bumps, Humidity, and Drying Glitches:} Bumpy tissue surfaces, tilted angles, room humidity, skin creams, or a lab sample drying out during scanning can severely skew terahertz readings.
\end{itemize}

\section{Dealing with Uncertainties}
\begin{itemize}
  \item \textbf{Shallow Water Barrier:} Heavy absorption by water stops terahertz waves after 0.1 to 2~millimeters, so deeper tumor spread or internal organs cannot be seen.
  \item \textbf{Hard to Separate Benign Hydration from Cancer:} Dense breast glands, scars, and normal swelling overlap strongly with malignant terahertz readings.
  \item \textbf{Humidity, Tilt, and Drying Errors:} Uneven surfaces, room moisture, creams, and drying surgical specimens can distort wave strength and timing.
  \item \textbf{No Proven Clinical Cutoffs Yet:} Because no terahertz system is clinically approved, research cutoffs and AI scores are not yet standardized across hospitals; building large, shared clinical terahertz tissue databases is an active step needed to make results repeatable.
  \item \textbf{Standard Care Always Decides:} Researchers cut down glitches with dry-nitrogen purging, mirror calibration, gentle blotting, and fast scans, while every real medical decision relies strictly on official lab pathology, dermoscopy, eye OCT, or standard breast imaging.
\end{itemize}

\section{Augmenting Methods}
\begin{itemize}
  \item \textbf{Official Lab Pathology (H\&E):} Standard pathology on removed breast, skin, or digestive samples remains the mandatory gold standard for checking tumor edges and deciding on further treatment.
  \item \textbf{Eye OCT and Skin Dermoscopy:} High-resolution OCT, corneal thickness checks (pachymetry), and dermoscopy give sharp cellular depth views alongside terahertz water readings.
  \item \textbf{Mammography, Specimen X-Rays, and Ultrasound:} Standard breast scans and operating-room specimen X-rays spot tiny calcium specks and deep borders that surface terahertz waves cannot reach.
  \item \textbf{Confocal Microscopy and Fluorescence:} Researchers pair terahertz hydration maps with optical microscopy or glowing dyes to help separate cancer from dense benign tissue.
  \item \textbf{Machine Learning Software:} Computer classifiers and water-behavior models help separate overlapping tumor and dense-tissue signals.
\end{itemize}

\section{Comparison with Alternative Modalities}
\begin{itemize}
  \item \textbf{Versus Optical Scans (OCT / Confocal):} OCT and confocal microscopy are 10 to 100 times sharper and are already proven everyday tools in eye and skin clinics, while terahertz imaging uniquely highlights surface water balance at coarser detail.
  \item \textbf{Versus Raman Imaging:} Raman ignores water and reads specific protein and fat chemical bonds, while terahertz is dominated by water---limiting depth, but making it sensitive to swelling and burns.
  \item \textbf{Versus Microwave Imaging:} Terahertz gives sharper surface detail than centimeter-scale microwaves, while microwaves reach several centimeters deep where terahertz waves cannot go.
  \item \textbf{Versus Lab Pathology and Specimen X-Rays:} Terahertz is strictly a research tool, whereas specimen X-rays and formalin-fixed H\&E pathology remain the required clinical standards for surgical margins.
\end{itemize}

\section{Key References}
\begin{itemize}
  \item Ashworth PC, Pickwell-MacPherson E, Provenzano E, et al. Terahertz pulsed spectroscopy of freshly excised human breast cancer. \textit{Optics Express}. 2009;17(15):12444--12454.
  \item Wallace VP, Fitzgerald AJ, Shankar S, et al. Terahertz pulsed imaging of basal cell carcinoma ex vivo and in vivo. \textit{British Journal of Dermatology}. 2004;151(2):424--432.
  \item Yang X, Zhao X, Yang K, et al. Biomedical applications of terahertz spectroscopy and imaging. \textit{Trends in Biotechnology}. 2016;34(10):810--824.
  \item International Commission on Non-Ionizing Radiation Protection (ICNIRP). Guidelines for limiting exposure to electromagnetic fields (100 kHz to 300 GHz). \textit{Health Physics}. 2020;118(5):483--524.
  \item Cong M, Li W, Liu Y, et al. \textit{Biomedical application of terahertz imaging technology: a narrative review}. Quantitative Imaging in Medicine and Surgery (NIH PubMed Central). Available at: \url{https://pmc.ncbi.nlm.nih.gov/articles/PMC10722018/}.
\end{itemize}

\section{Further Reading}
\begin{itemize}
  \item Wikipedia contributors. \textit{Terahertz radiation}. Wikipedia, The Free Encyclopedia. Available at: \url{https://en.wikipedia.org/wiki/Terahertz_radiation}.
\end{itemize}
